\maketitle

\begin{abstract}
Stone's representation theorem identifies a Boolean algebra $\E$ with the clopen
algebra of a compact totally-disconnected space $S(\E)$, and $\sigma(\E)$ with the
Baire $\sigma$-algebra of $S(\E)$.  Under this correspondence, the classical question
--- when does a charge on $\E$ extend to a $\sigma$-additive measure on $\sigma(\E)$?
--- becomes: when does a finitely-additive measure on the clopens of $S(\E)$ extend
to a Borel measure?  The answer for a single algebra is well known: continuity at
$\emptyset$, equivalently tightness on the Stone space.

This paper asks the question for a \emph{directed system} of Boolean algebras: a
query system, in which each algebra $\E_Q$ represents one observer's event language,
connected by refinement maps with compatible marginals.  The Stone spaces form a
projective system, and the extension question becomes whether charges propagate
coherently to the Baire $\sigma$-algebras throughout the system.

We show three things.  First, the gap $\sigE \setminus \E$ is not a defect: an
observer with no gap admits no nontrivial refinement and has no horizon.
Incompleteness is constitutive of discriminability.  Second, index-layer coherence
conditions alone --- sequential upper-directedness and compatibility --- do not force
$\sigma$-additivity; this confirms, in the system setting, that no finitary structural
condition suffices, consistent with the classical single-algebra result.  The
counterexample is a compatible family on a sequentially upper-directed system in which
every observer permanently assigns unit mass to a sequence of events the system as a
whole sees as empty.  Third, we identify the correct characterisation: a family of
finitely-additive contents extends to $\sigma$-additive measures at every level if and
only if it is \emph{collectively exhaustive} --- a valuation-layer condition requiring
that every vanishing sequence is witnessed somewhere in the refinement order.
Collective exhaustion is the algebraic substitute for the tightness condition on the
Stone space, requiring no topology on the outcome spaces.
\end{abstract}

%%%%%%%%%%%%%%%%%%%%%%%%%%%%%%%%%%%%%%%%%%%%%%%%%%%%%%%%%%%%%%%%%%%%%
\section{Introduction}
%%%%%%%%%%%%%%%%%%%%%%%%%%%%%%%%%%%%%%%%%%%%%%%%%%%%%%%%%%%%%%%%%%%%%

A Boolean algebra $\E$ of observable events, together with a finitely-additive
valuation $\ell : \E \to [0,1]$, is the basic data of a discriminating observer.
Under Stone's representation theorem~\cite{Stone1936}, $\E$ is isomorphic to the
clopen algebra of a compact totally-disconnected space $S(\E)$, and $\sigma(\E)$ is
the Baire $\sigma$-algebra of $S(\E)$.  The extension question --- when does $\ell$
extend to a $\sigma$-additive measure on $\sigma(\E)$? --- becomes: when is the
charge on the clopens tight enough to extend to the full Baire $\sigma$-algebra?
The extension criterion (Theorem~\ref{thm:extension-criterion}) gives the answer:
continuity at $\emptyset$, in the sense of Daniell--Stone~\cite{Daniell1918,Stone1948},
is necessary and sufficient.  No finitary condition on $\E$ alone can supply this:
the Stone space of a non-compact Boolean algebra carries charges that do not extend,
regardless of any algebraic constraint~\cite{deFinetti1974}.  The question is
whether something beyond the single algebra can.

This paper studies the extension question for a \emph{directed system} of Boolean
algebras: a query system, in which each $\E_Q$ is one observer's event language,
connected by refinement maps and compatible marginals.  The Stone spaces $\{S(\E_Q)\}$
form a projective system.  The question becomes whether the charges propagate
coherently to $\sigma$-additive measures throughout --- and if so, what condition on
the family is necessary and sufficient.

The paper establishes three results.  The negative result (Section~\ref{sec:negative})
shows that the extension question has no finitary answer even in the system setting:
the structural regress is irreducible, and no chain of finite refinements resolves it.
The incompleteness result (Section~\ref{sec:incompleteness}) shows the gap
$\sigE \setminus \E$ is necessary: an observer with $\E_Q = \sigE_Q$ admits no
nontrivial refinement and cannot be embedded in any interesting query system.  The
positive result (Section~\ref{sec:positive}) identifies the correct characterisation.
The natural question --- whether index-layer conditions (directedness, compatibility)
force $\sigma$-additivity --- turns out to be malformed: it conflates conditions on
the query system's structure with conditions on the contents it carries.  We show that
index-layer coherence does not force $\sigma$-additivity (the finite-cofinite
counterexample), and that $\sigma$-additivity is equivalent to collective exhaustion,
a purely valuation-layer condition (Theorem~\ref{thm:sp1}).

Section~\ref{sec:discussion} locates the results and states what remains open.

\subsection{Organisation}

Section~\ref{sec:setup} sets up the single-observer case: a Boolean algebra $\E$ of
distinctions and a finitely-additive valuation $\ell$ on it, with no topology and no
$\sigma$-algebra assumed.

Section~\ref{sec:negative} establishes the negative result: exhaustiveness is trivial,
the real obstruction is continuity at $\emptyset$ for sequences escaping $\E$, and no
finitary condition on any single level can supply this.

Section~\ref{sec:incompleteness} establishes the necessity of incompleteness: an
observer with $\E_Q = \sigE_Q$ admits no genuine refinement.

Section~\ref{sec:positive} completes the analysis: it proves the independence of
index- and valuation-layer conditions and establishes the main equivalence (collective
exhaustion $\iff$ $\sigma$-additive extensibility).

Section~\ref{sec:discussion} locates the results within the Stone duality picture and
states what remains open.

%%%%%%%%%%%%%%%%%%%%%%%%%%%%%%%%%%%%%%%%%%%%%%%%%%%%%%%%%%%%%%%%%%%%%
\section{The single-observer setup}
\label{sec:setup}
%%%%%%%%%%%%%%%%%%%%%%%%%%%%%%%%%%%%%%%%%%%%%%%%%%%%%%%%%%%%%%%%%%%%%

We begin with a single Boolean algebra and its charge, setting aside the query system
structure to isolate the extension question.

\begin{definition}[Discriminability structure]
A \emph{discriminability structure} is a pair $(\E, \ell)$ where:
\begin{enumerate}[label=(\roman*)]
  \item $\E$ is a Boolean algebra of subsets of some set $O$ --- closed under finite
    union, intersection, and complementation, containing $\emptyset$ and $O$.
  \item $\ell : \E \to [0,1]$ is a normalized finitely-additive valuation:
    $\ell(\emptyset) = 0$, $\ell(O) = 1$, and $\ell(A \cup B) = \ell(A) + \ell(B)$
    whenever $A \cap B = \emptyset$.
\end{enumerate}
Via Stone's representation theorem, $\E$ is isomorphic to the clopen algebra of
a compact totally-disconnected space $S(\E)$, the Stone space of $\E$, and
$\sigma(\E)$ is the Baire $\sigma$-algebra of $S(\E)$.
\end{definition}

\noindent No $\sigma$-algebra is assumed on $O$.  No topology is assumed.

\begin{definition}[The gap]
\label{def:gap}
For a discriminability structure $(\E, \ell)$, the \emph{gap} is the set
$\sigE \setminus \E$ --- events in the $\sigma$-algebra generated by $\E$ that are
not themselves in $\E$.
\end{definition}

\noindent The gap consists of events reachable by countable Boolean operations on
$\E$ but not contained in $\E$ itself.

\begin{definition}[Extension]
We say $\ell$ \emph{extends} to $\sigE$ if there exists a $\sigma$-additive measure
$\mu$ on $\sigE$ with $\mu|_\E = \ell$.  The extension is \emph{unique} if any two
such $\mu$ agree on all of $\sigE$.
\end{definition}

\noindent The central question of this paper is: what properties of $(\E, \ell)$ are
necessary and sufficient for such an extension to exist and be unique?

%%%%%%%%%%%%%%%%%%%%%%%%%%%%%%%%%%%%%%%%%%%%%%%%%%%%%%%%%%%%%%%%%%%%%
\section{The negative result: finitary conditions cannot suffice}
\label{sec:negative}
%%%%%%%%%%%%%%%%%%%%%%%%%%%%%%%%%%%%%%%%%%%%%%%%%%%%%%%%%%%%%%%%%%%%%

\subsection{Exhaustiveness is trivial}

The Stone space $S(\E)$ of a Boolean algebra $\E$ is compact by construction.
But compactness of $S(\E)$ does not force a given charge on $\E$ to extend to a
Borel measure on $S(\E)$: compactness is a property of the space, not the measure.
The measure-theoretic condition is tightness --- equivalently, continuity at
$\emptyset$.  The question is whether this can be witnessed finitely.  It cannot.

The first natural candidate is exhaustiveness.

\begin{definition}[Exhaustiveness]
A finitely-additive valuation $\ell : \E \to [0,1]$ is \emph{exhaustive} if
$\ell(E_n) \to 0$ for every disjoint sequence $(E_n)_{n \geq 1} \subset \E$.
\end{definition}

\begin{proposition}
\label{prop:exhaustiveness-trivial}
Every normalized finitely-additive valuation on a Boolean algebra is exhaustive.
\end{proposition}

\begin{proof}
Let $(E_n)$ be a disjoint sequence in $\E$.  For any $N$,
$\sum_{n=1}^N \ell(E_n) = \ell\bigl(\bigcup_{n=1}^N E_n\bigr) \leq \ell(O) = 1$
by finite additivity and normalization.  Hence $\sum_{n=1}^\infty \ell(E_n) \leq 1$,
so $\ell(E_n) \to 0$.
\end{proof}

\noindent Exhaustiveness carries no information about extensibility.  It is a
consequence of normalization alone, independent of any coherence condition or
structural hypothesis.

\subsection{The real obstruction}

The correct extension condition is continuity from above at $\emptyset$.

\begin{definition}[Continuity at $\emptyset$]
\label{def:cont-empty}
A finitely-additive valuation $\ell : \E \to [0,1]$ is \emph{continuous at
$\emptyset$} if: for every decreasing sequence $E_1 \supseteq E_2 \supseteq \cdots$
in $\E$ with $\bigcap_{n=1}^\infty E_n = \emptyset$ in $\sigE$, we have $\ell(E_n)
\to 0$.
\end{definition}

\begin{remark}
The intersection $\bigcap_n E_n$ in Definition~\ref{def:cont-empty} is taken in
$\sigE$, not in $\E$.  The condition is therefore not verifiable from within $\E$
alone: confirming $\bigcap_n E_n = \emptyset$ requires access to $\sigE$.  This
is precisely what the regress below captures.
\end{remark}

\begin{theorem}[Extension criterion]
\label{thm:extension-criterion}
A normalized finitely-additive $\ell : \E \to [0,1]$ extends to a $\sigma$-additive
measure on $\sigE$ if and only if $\ell$ is continuous at $\emptyset$.  The extension
is then unique.
\end{theorem}

\begin{proof}
Standard: see Halmos~\cite{Halmos1950} or Fremlin~\cite{Fremlin2003} Vol.~1.
The forward direction is immediate from $\sigma$-additivity of the extension.
For the converse, continuity at $\emptyset$ is precisely the condition used in the
standard proof of the Carath\'{e}odory extension to pass from finite additivity to
$\sigma$-additivity on the generated $\sigma$-algebra.
\end{proof}

\subsection{The regress}

In a query system, one might hope that a finer query $Q' \succeq Q$ can witness
the emptiness that $Q$ cannot see, and that compatibility then forces continuity at
$\emptyset$ at level $Q$.

\begin{proposition}[The regress]
\label{prop:regress}
Let $Q' \succeq Q$ with refinement map $\pi : O_{Q'} \to O_Q$.  Suppose
$E_1 \supseteq E_2 \supseteq \cdots$ in $\E_Q$ with $\bigcap_n \pi^{-1}(E_n) =
\emptyset$ in $\E_{Q'}$.  Compatibility gives $\ell_Q(E_n) = \ell_{Q'}(\pi^{-1}(E_n))$.
But this does not force $\ell_{Q'}(\pi^{-1}(E_n)) \to 0$: the same obstruction
reappears at level $Q'$, since finite additivity on $\E_{Q'}$ alone does not imply
continuity at $\emptyset$ for sequences decreasing to an element of $\E_{Q'}$.
\end{proposition}

\noindent The regress is structural.  Pushing the witnessing one level up does not
resolve it: it only relocates it.  No finite chain of refinements can supply continuity
at $\emptyset$ from purely finitary conditions.

\begin{corollary}
\label{cor:no-finitary}
No condition stateable purely in terms of the Boolean algebra structure of $\E_Q$ and
the finitely-additive valuations $\{\ell_{Q'}\}_{Q' \succeq Q}$ at finitely many
levels is necessary and sufficient for $\ell_Q$ to extend to $\sigE_Q$.
\end{corollary}

\noindent Something external to any single level is required.  The three candidates
are: topology (compactness witnesses emptiness), a limit query (a colimit in the
directed system whose event algebra contains the limit events), or an explicit
continuity axiom (honest but circular).  The remainder of this paper develops the
limit query option, which is the most intrinsic to the observer-system framework.

%%%%%%%%%%%%%%%%%%%%%%%%%%%%%%%%%%%%%%%%%%%%%%%%%%%%%%%%%%%%%%%%%%%%%
\section{The necessity of incompleteness}
\label{sec:incompleteness}
%%%%%%%%%%%%%%%%%%%%%%%%%%%%%%%%%%%%%%%%%%%%%%%%%%%%%%%%%%%%%%%%%%%%%

Before turning to the positive result, we establish that the gap is not a defect
to be eliminated.

\begin{definition}[Nontrivial refinement]
\label{def:nontrivial}
A refinement $Q' \succ Q$ (strict) is \emph{nontrivial} if there exists
$A' \in \E_{Q'}$ such that $\pi^{-1}(\pi(A')) \neq A'$ --- that is, the refinement
map does not preserve all distinctions in $\E_{Q'}$ back to $\E_Q$.  Equivalently,
$Q'$ makes distinctions that $Q$ cannot.
\end{definition}

\begin{theorem}[Discriminability requires incompleteness]
\label{thm:incompleteness}
Let $Q$ be a query admitting at least one nontrivial refinement $Q' \succ Q$.
Then $\E_Q \subsetneq \sigE_Q$: the gap is nonempty.
\end{theorem}

\begin{proof}
Since $Q' \succ Q$ is nontrivial, there exists $A' \in \E_{Q'}$ with
$\pi^{-1}(\pi(A')) \neq A'$.  Let $D = A' \setminus \pi^{-1}(\pi(A'))$ --- a
nonempty set in $\E_{Q'}$ that maps to the same coarse event as part of $A'$ but
is genuinely finer.  The preimage $\pi^{-1}(B)$ for any $B \in \E_Q$ is a union
of level-$Q'$ atoms, hence cannot separate $D$ from its complement within
$\pi^{-1}(\pi(A'))$.  It follows that the event $\pi^{-1}(\pi(A'))$ in $\E_Q$,
when pulled back, conflates distinctions that $\E_{Q'}$ separates.  The sequence
of such distinctions generates events in $\sigE_Q$ --- specifically, the intersection
of the coarser level events --- that are not in $\E_Q$.
\end{proof}

\begin{remark}
The contrapositive of Theorem~\ref{thm:incompleteness} is equally important: a
query $Q$ with $\E_Q = \sigE_Q$ cannot be nontrivially refined.  Any $Q' \succeq Q$
satisfies $\E_{Q'} = \pi^{-1}(\E_Q)$, so the refinement order above $Q$ collapses.
The gap $\sigE_Q \setminus \E_Q$ is therefore not a deficiency of the construction
but a necessary feature of any query that admits genuine refinement.
\end{remark}

\begin{corollary}
In any query system with a genuine (nontrivial, infinite) refinement order, every
observer at a finite level $Q$ has a nonempty gap.  Discriminability is necessarily
incomplete at every finite level.
\end{corollary}

%%%%%%%%%%%%%%%%%%%%%%%%%%%%%%%%%%%%%%%%%%%%%%%%%%%%%%%%%%%%%%%%%%%%%
\section{Collective exhaustion and the characterisation of $\sigma$-additivity}
\label{sec:positive}
%%%%%%%%%%%%%%%%%%%%%%%%%%%%%%%%%%%%%%%%%%%%%%%%%%%%%%%%%%%%%%%%%%%%%

The negative result established that no finitary or index-layer condition can force
$\sigma$-additivity.  This section completes the picture: we identify the exact
valuation-layer condition that characterises $\sigma$-additive extensibility, show it
is independent of the index-layer conditions, and prove the equivalence.

\begin{proposition}[Independence of index and valuation layers]
\label{prop:independence}
Sequential upper-directedness together with compatibility of finitely-additive contents
does not force $\sigma$-additivity.  There exists a compatible family of normalized
finitely-additive contents on a sequentially upper-directed query system that fails
$\sigma$-additive extension at every level.
\end{proposition}

\begin{proof}
Let $\iota = \mathbb{N} \cup \{\omega\}$ with $n \leq \omega$ for all $n \in \mathbb{N}$,
directed by this order.  This is sequentially upper-directed: $\omega$ is a universal
upper bound for every sequence.  Take all outcome spaces to be $\mathbb{Q}$, all event
algebras to be the finite-cofinite algebra $\mathcal{E} = \{E \subseteq \mathbb{Q} :
E \text{ is finite or } E^c \text{ is finite}\}$, and all refinement maps to be the
identity.  Define $\ell(E) = 0$ if $E$ is finite and $\ell(E) = 1$ if $E$ is cofinite.
This is a normalized finitely-additive content, and compatibility holds trivially.

Fix an enumeration $q : \mathbb{N} \to \mathbb{Q}$.  Let $F_n = \{q(0), \ldots, q(n)\}$.
Each $F_n$ is finite, so $\ell(F_n) = 0$ for all $n$.  Since $q$ is surjective,
$\bigcup_n F_n = \mathbb{Q}$, and $\ell(\mathbb{Q}) = 1$.  $\sigma$-subadditivity
would give $1 \leq \sum_n \ell(F_n) = 0$: contradiction.  So $\ell$ fails
$\sigma$-additive extension at every level.
\end{proof}

\noindent The counterexample establishes that index-layer coherence is insufficient:
the $\sigma$-algebra is not derivable from structural conditions on the query system
alone.  The content is perfectly coherent within the finite-cofinite algebra; the
failure is a valuation-layer pathology invisible to the index structure.

\begin{definition}[Collectively exhaustive]
\label{def:collective-exhaustion}
A compatible family $\{\ell_Q\}$ of normalized finitely-additive contents on a query
system is \emph{collectively exhaustive} if: for every level $Q$ and every decreasing
sequence $E_1 \supseteq E_2 \supseteq \cdots$ of events in $\E_Q$ with $\bigcap_n E_n
= \emptyset$, there exists $Q' \succeq Q$ such that $\ell_{Q'}(\pi^{-1}(E_n)) \to 0$.
\end{definition}

\noindent Collective exhaustion is a \emph{valuation-layer} condition: it places a
requirement on how the contents behave, not on the index structure of the query system.
It says that no observer can assign persistent mass to events that the system as a whole
sees as empty.  The counterexample above fails collective exhaustion: for the sequence
$E_n = \mathbb{Q} \setminus F_n$ (cofinite sets shrinking to $\emptyset$), every level
assigns content $1$ to $\pi^{-1}(E_n) = E_n$, and no witness exists anywhere in the
system.

\begin{theorem}[Collective exhaustion characterises $\sigma$-additivity]
\label{thm:sp1}
Let $\{\ell_Q\}$ be a normalized compatible family of finitely-additive contents on
a query system.  The following are equivalent:
\begin{enumerate}[label=(\roman*)]
  \item The family is collectively exhaustive.
  \item $\ell_Q$ extends uniquely to a $\sigma$-additive measure on $\sigE_Q$
        for every $Q$.
\end{enumerate}
\end{theorem}

\begin{proof}
\emph{(i) $\Rightarrow$ (ii).}  Fix $Q$ and a decreasing sequence $E_n \searrow
\emptyset$ in $\mathcal{E}_Q$.  By collective exhaustion, find $Q' \succeq Q$ with
$\ell_{Q'}(\pi^{-1}(E_n)) \to 0$.  By compatibility, $\ell_Q(E_n) =
\ell_{Q'}(\pi^{-1}(E_n)) \to 0$.  So $\ell_Q$ is continuous at $\emptyset$.  By the
extension criterion (Theorem~\ref{thm:extension-criterion}), $\ell_Q$ extends uniquely
to a $\sigma$-additive measure on $\sigE_Q$.

\emph{(ii) $\Rightarrow$ (i).}  Suppose $\ell_Q$ extends to $\mu_Q$ for every $Q$.
Let $E_n \searrow \emptyset$ in $\mathcal{E}_Q$.  Then $\ell_Q(E_n) = \mu_Q(E_n) \to
\mu_Q(\emptyset) = 0$ by $\sigma$-additivity of $\mu_Q$ and normalization (which ensures
$\mu_Q(E_0) \leq \mu_Q(O_Q) = 1 < \infty$, the finiteness condition for continuity from
above).  Take $Q' = Q$.
\end{proof}

\begin{remark}[The dissolution]
\label{rem:dissolution}
The original witnessing conjecture asked whether index-layer conditions (sequential
upper-directedness, compatibility) force $\sigma$-additivity.
Proposition~\ref{prop:independence} shows the answer is no.  The question was
malformed: it conflated the index layer (the preorder, refinement maps, sequential
upper-directedness) with the valuation layer (the contents, compatibility, continuity
at $\emptyset$).  These layers are independent.

The correct question is: what valuation-layer condition characterises $\sigma$-additive
extensibility?  Theorem~\ref{thm:sp1} answers it: collective exhaustion.  It is not an
axiom imposed from outside --- it is the exact semantic criterion separating systems that
model coherent worlds (every convergence is witnessed somewhere in the refinement order)
from those that do not.

Note that sequential upper-directedness plays no role in Theorem~\ref{thm:sp1}: it is
a per-level result, mediated by the refinement structure.  Sequential upper-directedness
enters when one asks whether the per-level extensions assemble into a global probability
space over $\Omega$ --- that is the next question.
\end{remark}

%%%%%%%%%%%%%%%%%%%%%%%%%%%%%%%%%%%%%%%%%%%%%%%%%%%%%%%%%%%%%%%%%%%%%
\section{Discussion}
\label{sec:discussion}
\label{sec:discussion}
%%%%%%%%%%%%%%%%%%%%%%%%%%%%%%%%%%%%%%%%%%%%%%%%%%%%%%%%%%%%%%%%%%%%%

Under Stone duality, $\E$ corresponds to the clopen algebra of a compact
totally-disconnected space and $\sigma(\E)$ to its Baire $\sigma$-algebra; the
extension question is whether a charge on the clopens extends to a Borel measure.
The present paper answers this for directed systems without topological hypotheses
on the outcome spaces: collective exhaustion is the necessary and sufficient
condition, playing the role that tightness plays on the Stone space.

Theorem~\ref{thm:sp1} is a per-level result.  Sequential upper-directedness plays
no role in it: that condition enters when asking whether the per-level extensions
$\{\mu_Q\}$ assemble into a single $\sigma$-additive measure on the projective
limit $\Omega = \varprojlim O_Q$.  For compact outcome spaces this follows from
the Prokhorov extension theorem; for standard Borel spaces from Musia\l's theorem
via perfectness.  The topology-free case remains open.

The incompleteness theorem (Theorem~\ref{thm:incompleteness}) establishes that the
gap $\sigE_Q \setminus \E_Q$ is nonempty whenever $Q$ admits a nontrivial
refinement.  Making the proof fully explicit --- constructing a specific element
of $\sigE_Q \setminus \E_Q$ from the refinement data --- connects to the Stone
space picture: $\sigE_Q$ is the $\sigma$-closure of $\E_Q$ in the Baire
$\sigma$-algebra of $S(\E_Q)$, and nontrivial refinements correspond to points
of $S(\E_Q)$ not isolated in the clopen topology.

\bibliographystyle{plain}
\bibliography{references}

\section*{Acknowledgements}

The author thanks Claude (Anthropic) for assistance with mathematical writing and
technical discussion throughout the development of this work.  The author gratefully
acknowledges financial support from Dr.\ Simon Bonner and Dr.\ Matt Davison at the
University of Western Ontario, and from MITACS.
